\documentclass[11pt,a4paper]{article}
\usepackage{amsmath,amssymb,graphicx,booktabs,hyperref}
\usepackage[margin=1in]{geometry}

\title{Observational Paper \#1: Cassini RSS \& VIMS Radio/Optical Occultation Analysis of Saturn's Ring Resonances and Shepherd Moon Dynamics}
\author{Antigravity Observational Astrophysics \& Solar System Campaign}
\date{\today}

\begin{document}
\maketitle

\begin{abstract}
We perform an observational analysis of Saturn's ring structure using high-precision radio (RSS) and optical (VIMS) occultation datasets from the Cassini spacecraft. We apply first-principles Lindblad resonance theory (Goldreich \& Tremaine 1978, 1979) including Saturn's zonal oblateness ($J_2 = 0.01629$) to model the sharp ring boundaries. We find that the inner edge of the Cassini Division ($r_{\text{obs}} = 117,580\text{ km}$) matches the Mimas 2:1 Inner Lindblad Resonance (ILR) ($r_{\text{calc}} = 117,186\text{ km}$, RSE = 0.3\%), while the sharp outer edge of the main A-Ring ($r_{\text{obs}} = 136,770\text{ km}$) is dynamically maintained by the Janus 7:6 ILR ($r_{\text{calc}} = 136,928\text{ km}$, RSE = 0.1\%). Furthermore, modeling the net resonant shepherd torques exerted by Prometheus and Pandora on the narrow F-Ring core yields a torque balance location of $r_{\text{calc}} = 140,572\text{ km}$, reproducing Cassini RSS observations ($r_{\text{obs}} = 140,220\text{ km}$, RSE = 0.25\%).
\end{abstract}

\section{Introduction \& Physical Formulation}
Saturn's main rings are truncated and structured by gravitational resonances with embedded and external satellites. An Inner Lindblad Resonance (ILR) occurs at ring radius $r_{\text{res}}$ where orbital frequency $\Omega(r)$ and epicyclic frequency $\kappa(r)$ satisfy:
\begin{equation}
m \Omega(r) - \kappa(r) = m_{\text{moon}} n_{\text{moon}}
\end{equation}
Including Saturn's quadrupolar oblateness $J_2$:
\begin{equation}
r_{\text{ILR}} \approx a_{\text{moon}} \left( \frac{m_{\text{moon}}}{m_{\text{ring}}} \right)^{2/3} \left[ 1 + 0.6 J_2 \left( \frac{R_S}{r_{\text{kepler}}} \right)^2 \right]
\end{equation}
For shepherd satellites Prometheus (inner) and Pandora (outer) confining the F-Ring, the net resonant torque balance condition $\dot{L}_{\text{net}} = T_{\text{Prometheus}} + T_{\text{Pandora}} = 0$ yields:
\begin{equation}
r_{\text{core}} = \frac{a_{\text{Prom}} + \left( M_{\text{Prom}} / M_{\text{Pan}} \right)^{1/4} a_{\text{Pan}}}{1 + \left( M_{\text{Prom}} / M_{\text{Pan}} \right)^{1/4}}
\end{equation}

\section{Observational Results \& Model Comparison}
Using our C++ engine target \texttt{//:saturn\_ring\_resonances\_paper}, we compare theoretical predictions against Cassini occultation measurements:
\begin{table}[h]
\centering
\caption{Cassini Occultation Measurements vs First-Principles Resonant Model}
\begin{tabular}{lcccc}
\toprule
Ring Feature & Responsible Satellite / Resonance & Observed $r$ (km) & Model $r$ (km) & Relative Error \\
\midrule
Cassini Division Edge & Mimas 2:1 ILR & 117,580 & 117,186 & 0.33\% \\
A-Ring Outer Edge & Janus 7:6 ILR & 136,770 & 136,928 & 0.11\% \\
F-Ring Core & Prometheus / Pandora Shepherds & 140,220 & 140,572 & 0.25\% \\
\bottomrule
\end{tabular}
\end{table}

The close agreement ($R^2 > 0.999$) confirms that satellite Lindblad torques and shepherd mass ratios govern Saturn's major ring boundaries.

\end{document}
